\section{Mengapa Teknik Informatika Memuja Logika?}

Selamat datang di titik awal medan tempur Ilmu Komputer (*Computer Science*). Saat Anda membuka lembaran buku ini, Anda mungkin bertanya-tanya: \textit{"Untuk apa seorang mahasiswa Teknik Informatika yang seharusnya belajar mengetik kode aplikasi, malah dipaksa duduk berjibaku mempelajari deretan huruf P, Q, dan operator V (OR), $\land$ (AND) yang terlihat seperti peninggalan filsuf Yunani Kuno?"} \cite{rosen2019,copi2014}

Jawabannya keras dan sederhana: \textbf{Mengetik kode (Coding) bisa dilakukan oleh siapa saja, bahkan oleh anak-anak menggunakan *Drag-and-Drop*. Namun, Mendesain Arsitektur Rekayasa Perangkat Lunak (\textit{Software Engineering}) dan Algoritma Kecerdasan Buatan (\textit{Artificial Intelligence}) yang bebas \textit{Bug} maut, murni mensyaratkan otak yang telah dikalibrasi oleh Logika Matematika} \cite{levin_discrete,stanford_cs103}.

Buku ajar ini ditulis bukan sebagai seonggok dokumen matematis teoretis yang mengantuk. Ini adalah **Manual Keselamatan (\textit{Survival Guide})** Anda untuk diizinkan masuk ke dalam jajaran Insinyur IT berstandar global. Universitas terkemuka dunia seperti MIT (dalam \textit{Mathematics for Computer Science}) dan Stanford (dalam \textit{Mathematical Foundations of Computing}) \cite{mit_6042,stanford_cs103} secara absolut memelopori bahwa tanpa Logika Proposisional dan Predikat, seorang sarjana komputer tak lebih dari sekadar "Tukang Ketik Berijazah".

\subsection{Menyelamatkan Sistem dari Kiamat \textit{Logical Bug}}

Aplikasi kasir, sistem keamanan reaktor nuklir, hingga *Autopilot* pesawat Boeing, semuanya dijalankan oleh deretan jutaan *If-Else Statement* (Kondisi Percabangan). Ketika seorang *Programmer* buta akan hukum \textit{Kontraposisi} atau gagal mengadopsi \textit{De Morgan's Laws} saat meracik filter *Password* raksasa, mesin bukannya *Error* secara sintaks. Mesin akan terus berjalan mulus, **namun** diam-diam membelokkan arah pesawat menabrak gunung atau membocorkan data triliunan saldo Bank. Di ranah komputasi, kecacatan alur pikir tak kasatmata ini bergelar \textbf{\textit{Logic Bug}}---mimpi buruk yang hanya bisa diredam jika otak Anda terbiasa menguliti tabel kebenaran validitas Argumen \cite{forallx,iep_natded}.

\subsection{Tangga Menuju Kecerdasan Buatan (AI)}

Memasuki era abad 21, Anda niscaya bercita-cita membangun *Machine Learning* atau *Expert Systems* yang sanggup mendiagnosis kanker sehebat dokter \cite{saylor_cs202}. AI modern tidak menebak; ia bernalar murni menggunakan struktur \textbf{Logika Predikat Orde Pertama (\textit{First-Order Logic})}. Buku ini akan mengajarkan langkah demi langkah bagaimana menarik Kesimpulan Mutlak (\textit{Inference}) bertenaga *Modus Ponens*, mengawinkan insting manusia purba dengan komputasi biner abad mesin. Menguasai buku ini, berarti Anda menggenggam cetak biru (*blueprint*) sejati bagaimana *Database Relasional (SQL)* dan Mesin *AI* mengambil keputusan \cite{rosen2019}.
